\section{An alternative route: one loop for unordered proof certificates}
\label{sec:single-loop}

This section gives a source-level improvement independent of Turing-machine
simulation and Jones compression.  It does not claim a new measured final
polynomial.  The affine-product construction in package~68 uses one bounded
universal block for proof steps and a second for a common-multiple condition.
Padding the proof makes the two blocks share one index.  The resulting source
has four outer natural parameters instead of five and one bounded universal
block instead of two.  The three separate label bounds are merged into one joint bound, reducing
its local witness count from $23$ to $22$.
The number of witnesses after arithmetic closure has not been evaluated.

\subsection{Statement}

Let a Horn calculus have natural-number labels, a freely derivable truth
label~$0$, and rules $(p,q)\longrightarrow c$ with one conclusion and two
premise slots.  Axioms have premises $(0,0)$; unary rules are padded by~$0$.
Include the harmless rule $(0,0)\longrightarrow0$ and fix target~$1$.
Write $D$ for its least set of derivable labels.

For $M\geq1$, put $f_M(j)=1+jM$.  Let
\[
 \beta(x,y,i)=x\bmod\bigl(1+(i+1)y\bigr).
\]
The following predicate has outer natural parameters $T,x,y,M_0$, where
$M=M_0+1$:
\begin{align}
 \beta(x,y,0)&=M+1,& \beta(x,y,T)&=1, \label{eq:sl-endpoint}\\
 \forall i<T\quad\exists p,q,c,d,k\in\mathbb N:\quad
 (p,q)&\longrightarrow c, \nonumber\\
 p+q+c+d&=T,\label{eq:sl-bounds}\\
 M&=(i+1)k,\label{eq:sl-common}\\
 \beta(x,y,i+1)f_M(c)&=
 \beta(x,y,i)f_M(p)f_M(q).\label{eq:sl-transition}
\end{align}
The rule relation in this display can have its own fixed local witnesses.
Only one occurrence of bounded universal quantification is present.

\paragraph{Theorem.}
The predicate has a solution if and only if $1\in D$.
This equivalence is provable by elementary finite mathematics in ZFC.

\subsection{Soundness}

The endpoints imply $T>0$: at $T=0$ the same beta value would equal both
$M+1\geq2$ and~$1$.  Equation~\eqref{eq:sl-common}, as $i$ runs from~$0$
to~$T-1$, says that every positive integer at most~$T$ divides~$M$.
For distinct $1\leq a<b\leq T$, the weights $f_M(a)$ and $f_M(b)$ are
coprime.  Indeed, a common divisor $d$ divides $(b-a)M$ and is coprime
to~$M$, hence divides $b-a$.  But $b-a\mid M$, forcing $d=1$.

Put $S_i=\beta(x,y,i)$.  Since all the weights are positive, $S_0=M+1>0$
and~\eqref{eq:sl-transition} imply $S_i>0$ for every $i\leq T$.
Multiplying the transition identities and cancelling the positive intermediate
states gives
\[
 \prod_{i<T} f_M(c_i)
 =f_M(1)\prod_{i<T}f_M(p_i)f_M(q_i).
\]
Pairwise coprimality of the positive-label weights, together with $f_M(0)=1$,
forces exact equality of positive-label multiplicities.  For example, choose
one prime divisor of each $f_M(j)$ and compare its valuation; that divisor
occurs in no other weight.  In the free abelian group on positive labels,
\[
 \sum_{i<T}(e_{c_i}-e_{p_i}-e_{q_i})=e_1,
 \qquad e_0:=0.
\]
If $1\notin D$, sum the finitely many coordinates outside~$D$.  A rule with
conclusion outside~$D$ has at least one premise outside~$D$, while a rule
with conclusion inside~$D$ contributes at most zero.  Thus the left side
has nonpositive sum and the right side has sum~$1$, a contradiction.

This argument does not accept cyclic proofs.  Its use of one conclusion per
rule is essential; the corresponding assertion for several conclusions is
false.

\subsection{Completeness and the order of the choices}

Take a finite proof tree of label~$1$.  Expand it from its root, recording a
rule every time an open nonzero obligation is replaced by its two premises;
zero premises are discarded.  Let $L$ be the number of recorded rules and
$m$ the greatest sum $p+q+c$ in a recorded rule.  Choose
\[
 T\geq\max(1,L,m),\qquad M=T!,\qquad M_0=M-1.
\]
Pad the list to length~$T$ with $(0,0)\longrightarrow0$.  These padding steps
leave the accumulator unchanged because $f_M(0)=1$.

At each stage let $S_i$ be the product of $f_M(j)$ over the current multiset
of nonzero open obligations.  Removing the chosen conclusion and inserting
its premises proves~\eqref{eq:sl-transition}.  The initial value is $M+1$;
the final value is~$1$.  Every row-label sum is at most~$T$, supplying the joint slack; and the quotients
$k_i=M/(i+1)$ establish~\eqref{eq:sl-common}.

Only after these finite integers have been chosen do we choose a beta scale.
For instance, take
\[
 y=T!\bigl(1+\max_{0\leq i\leq T}S_i\bigr).
\]
Then the moduli $1+(i+1)y$, for $0\leq i\leq T$, are pairwise coprime:
their index differences are at most~$T$ and divide~$y$.  They also exceed
the required residues.  The finite Chinese remainder theorem supplies
$x$ with $\beta(x,y,i)=S_i$ simultaneously.  This completes the proof.
The scale~$y$ and the common multiple~$M$ are distinct witnesses; identifying
them would generally destroy completeness.  The arithmetic closure
compiler's own scales must likewise remain independent.

\subsection{Literal local-polynomial adapter}

Package~68 supplies a nonnegative local polynomial
$H(S,S',m,M_0;\mathbf u)$ with $23$ natural witnesses and degree~$50$.
It checks one padded numeric C222 rule, its three label bounds, and its
affine accumulator update.  Write $H_0$ for its Horn product plus accumulator-update square, removing
the three old bound squares and their slack coordinates.  Here $\mathbf u$
denotes the twenty retained witnesses and includes the three label coordinates
$p,q,v$.  With a fresh joint-bound witness $b_0$, define
\[
 J(S,S',T,M_0,i;\mathbf u,b_0,k)
 =H_0(S,S',M_0;\mathbf u)
  + (p+q+v+b_0-T)^2 +\bigl(M_0+1-(i+1)k\bigr)^2.
\]
The retained nonnegativity makes $J=0$ precisely the conjunction of the rule/update, joint-bound, and divisor
conditions.  This source has five exposed step coordinates and $22$ local
natural witnesses.  Its degree remains exactly~$50$: the new summands
have degrees at most~$4$ and cannot remove the degree-$50$ part of the Horn product.

One now applies the supplied bounded-forall constructor once to the single
family $J=0$, conjoins the beta endpoint conditions, and existentially binds
the four outer parameters.  This is a precise source recipe, but the full
ordinary polynomial and its total cost have not been materialized here.
The old C222 calibration remains an inherited obligation if this source is
used as a Con(ZFC) endpoint.  No result in this section fills the repository's
separate gap between an external theorem and a checked internal ZFC proof.

The executable controls test $988$ artificial Horn systems with up to three
rules on labels $0,1,2$.  They construct and verify a certificate for all
$185$ systems in that family where target~$1$ is derivable, including beta
encoding, and reject all $185$ corrupted divisor quotients.  A separate
counterfeit uses $M=1$, axiom~$3$, and rule $(1,0)\longrightarrow0$:
the accumulator $2,4,1,1$ passes the raw update but fails the divisibility
condition.  These tests support implementation checks; the unbounded theorem
is the proof above.
